\section{付加体の構成岩・形成過程}

\subsection{平成27年度 問34}
\label{q:h27-reg-applied-34}

\noindent\textbf{原典：} sources/H27/17\_applied\_science.pdf、PDF物理13頁。

\textbf{公式正答：} 1（single）。

\par\textbf{独立解答と公式正答に相違あり：} 独立解答 5。初回独立解答5、公式1。公式表107頁を実画像で再確認。1の表層崩壊を否定する排他的根拠は未確定。5の大陸側斜面と海山の関係も原典のまま保持。メランジュの深層崩壊事例を表層崩壊の不可能性へ読み替えない。Astraで原典と地質機構を重点査読。

\paragraph{問題の要旨} 付加体の構造・弱面と斜面安定性に関する記述を判定する。

\Needspace{4\baselineskip}\paragraph{確認した条件・記述（独自要約）}\mbox{}\par

\noindent\textbf{方向：} 最も不適切なもの。\par

\noindent\textbf{対象：} 一般的な岩石分類・形成過程。\par

\paragraph{問われる核心} 付加体の不均質・異方性・斜面崩壊と未解決の肢解釈

\paragraph{必要知識} 付加体の不均質・異方性・斜面崩壊と未解決の肢解釈。母材・運搬堆積・有機物起源、風化、弱面、水の作用、岩盤と土の物性。

\paragraph{考え方・正答を選ぶ理由（査読待ち初稿）} 公式正答は1であり、公式正答表107頁のIII–34の1を画像でも確認した。初回独立解答は5である。原典1の「表層崩壊」と5の「大陸側斜面」「海山」の限定を含め、公式1を支持する排他的理由は現在の一次資料照合では確定していない。メランジュに深層崩壊事例があることだけから1の表層崩壊を誤りとは導かない。相違を未解決として保持し、公式を独立の推論結果と偽らない。

\Needspace{4\baselineskip}\paragraph{選択肢の検討（査読待ち初稿）}\mbox{}\par

\noindent\textbf{1：公式では不適切／根拠未確定} 砂岩の岩塊を剥離性の泥質基質が囲むメランジュは表層崩壊等を生じやすい。。公式はこの肢を不適切とする。しかしメランジュの硬軟混在と割れ目が切土等で問題となることは学会資料で確認できる。深層崩壊との関連もあるが、深層の事例だけで表層崩壊を不可能とは言えない。誤りとする限定根拠は未確定。\par

\noindent\textbf{2：概ね適切} 片理による結晶片岩の力学的異方性は地盤安定性に影響する。。面に沿うせん断抵抗などが方向で変わり、弱面と斜面の位置関係を評価する必要がある。等方な単一強度では表せない。\par

\noindent\textbf{3：条件付きで適切} 蛇紋岩には膨張性地山や光沢ある鱗片状剥離面を持ち、地すべりを起こしやすいものがある。。蛇紋岩地域の地すべりや粘土化した蛇紋岩の膨張性地山は一次資料に事例がある。健全な全蛇紋岩が膨潤するとの主張にはしない。\par

\noindent\textbf{4：概ね適切} 海溝の堆積物は圧縮で形成された衝上断層に沿って大陸側へ付加される。。付加は海洋プレート上の堆積物等が剥ぎ取られ、陸側に取り込まれる過程で、衝上断層が発達する。侵食型の沈み込み域もある。\par

\noindent\textbf{5：初回独立判断で疑義／未確定} 海溝の大陸側斜面で海底地すべりが起こり、海山が断裂して海溝へ崩壊する。。海底斜面崩壊自体は起こり得る。一方、海山は一般に沈み込む海洋プレート上の地形で、この一文の大陸側斜面との関係を疑って初回独立解答5とした。沈み込み中の海山破壊も起こり得るため、単に海山という語だけで絶対的な誤りとは確定しない。\par

\paragraph{再利用する考え方} 付加体の不均質・異方性・斜面崩壊と未解決の肢解釈を成因と観察される性質に結び付け、例示と全称命題、主要鉱物と副成分、公式解答と未解決の解釈を区別する。

\paragraph{共通理論} \hyperref[theory:engineering:accretion]{第\ref*{theory:engineering:accretion}節}, \hyperref[theory:engineering:anisotropy]{第\ref*{theory:engineering:anisotropy}節}

\par 参考：\cite{jseg_engineering_faq,kyoto_deep2018,sabo_deep2014,jsce_swelling2010,aist_accretion2026}。

\par 難易度 3（校準初稿）、優先度未評価。

\section{材料の応力・変形}

\subsection{平成29年度 問22}
\label{q:h29-reg-applied-22}

\noindent\textbf{原典：} sources/H29/17\_applied\_science.pdf、PDF物理7頁。

\textbf{公式正答：} 2（single）。

\par\textit{原典本文の校正と独自解説は未作成。分類のみの登録項目。}

\section{鉱物・岩石}

\subsection{平成24年度 問28}
\label{q:h24-reg-applied-28}

\noindent\textbf{原典：} sources/H24/17\_applied\_science.pdf、PDF物理8頁。

\textbf{公式正答：} 2（single）。

\par\textit{原典本文の校正と独自解説は未作成。分類のみの登録項目。}

\subsection{平成25年度 問28}
\label{q:h25-reg-applied-28}

\noindent\textbf{原典：} sources/H25/17\_applied\_science.pdf、PDF物理10頁。

\textbf{公式正答：} 2（single）。

\par\textit{原典本文の校正と独自解説は未作成。分類のみの登録項目。}

\section{堆積相・構造}

\subsection{平成24年度 問30}
\label{q:h24-reg-applied-30}

\noindent\textbf{原典：} sources/H24/17\_applied\_science.pdf、PDF物理9頁。

\textbf{公式正答：} 5（single）。

\par\textit{原典本文の校正と独自解説は未作成。分類のみの登録項目。}

\section{地盤・斜面災害}

\subsection{平成23年度 問33}
\label{q:h23-reg-applied-33}

\noindent\textbf{原典：} sources/H23/17\_applied\_science.pdf、PDF物理11頁。

\textbf{公式正答：} 5（single）。

\par\textit{原典本文の校正と独自解説は未作成。分類のみの登録項目。}

\subsection{平成23年度 問34}
\label{q:h23-reg-applied-34}

\noindent\textbf{原典：} sources/H23/17\_applied\_science.pdf、PDF物理11頁。

\textbf{公式正答：} 3（single）。

\par\textit{原典本文の校正と独自解説は未作成。分類のみの登録項目。}

\subsection{平成24年度 問31}
\label{q:h24-reg-applied-31}

\noindent\textbf{原典：} sources/H24/17\_applied\_science.pdf、PDF物理9頁。

\textbf{公式正答：} 5（single）。

\par\textit{原典本文の校正と独自解説は未作成。分類のみの登録項目。}

\subsection{平成25年度 問30}
\label{q:h25-reg-applied-30}

\noindent\textbf{原典：} sources/H25/17\_applied\_science.pdf、PDF物理11頁。

\textbf{公式正答：} 3（single）。

\par\textit{原典本文の校正と独自解説は未作成。分類のみの登録項目。}

\subsection{平成26年度 問33}
\label{q:h26-reg-applied-33}

\noindent\textbf{原典：} sources/H26/17\_applied\_science.pdf、PDF物理13頁。

\textbf{公式正答：} 2（single）。

\par\textit{原典本文の校正と独自解説は未作成。分類のみの登録項目。}

\subsection{平成26年度 問34}
\label{q:h26-reg-applied-34}

\noindent\textbf{原典：} sources/H26/17\_applied\_science.pdf、PDF物理14頁。

\textbf{公式正答：} 1（single）。

\par\textit{原典本文の校正と独自解説は未作成。分類のみの登録項目。}

\subsection{平成27年度 問35}
\label{q:h27-reg-applied-35}

\noindent\textbf{原典：} sources/H27/17\_applied\_science.pdf、PDF物理14頁。

\textbf{公式正答：} 5（single）。

\par\textit{原典本文の校正と独自解説は未作成。分類のみの登録項目。}

\subsection{平成28年度 問33}
\label{q:h28-reg-applied-33}

\noindent\textbf{原典：} sources/H28/17\_applied\_science.pdf、PDF物理13頁。

\textbf{公式正答：} 1（single）。

\par\textit{原典本文の校正と独自解説は未作成。分類のみの登録項目。}

\subsection{平成30年度 問24}
\label{q:h30-reg-applied-24}

\noindent\textbf{原典：} sources/H30/17\_applied\_science.pdf、PDF物理8頁。

\textbf{公式正答：} 3（single）。

\par\textit{原典本文の校正と独自解説は未作成。分類のみの登録項目。}

\subsection{平成30年度 問27}
\label{q:h30-reg-applied-27}

\noindent\textbf{原典：} sources/H30/17\_applied\_science.pdf、PDF物理9頁。

\textbf{公式正答：} 1（single）。

\par\textit{原典本文の校正と独自解説は未作成。分類のみの登録項目。}

\subsection{令和元年度再試験 問26}
\label{q:r01-ret-applied-26}

\noindent\textbf{原典：} sources/R01\_retest/17\_applied\_science.pdf、PDF物理9頁。

\textbf{公式正答：} 2（single）。

\par\textit{原典本文の校正と独自解説は未作成。分類のみの登録項目。}

\subsection{令和3年度 問28}
\label{q:r03-reg-applied-28}

\noindent\textbf{原典：} sources/R03/17\_applied\_science.pdf、PDF物理10頁。

\textbf{公式正答：} 2（single）。

\par\textit{原典本文の校正と独自解説は未作成。分類のみの登録項目。}

\subsection{令和3年度 問29}
\label{q:r03-reg-applied-29}

\noindent\textbf{原典：} sources/R03/17\_applied\_science.pdf、PDF物理10頁。

\textbf{公式正答：} 2（single）。

\par\textit{原典本文の校正と独自解説は未作成。分類のみの登録項目。}

\subsection{令和4年度 問27}
\label{q:r04-reg-applied-27}

\noindent\textbf{原典：} sources/R04/17\_applied\_science.pdf、PDF物理9頁。

\textbf{公式正答：} 4（single）。

\par\textit{原典本文の校正と独自解説は未作成。分類のみの登録項目。}

\section{岩石風化・変質・スレーキング}

\subsection{平成29年度 問24}
\label{q:h29-reg-applied-24}

\noindent\textbf{原典：} sources/H29/17\_applied\_science.pdf、PDF物理8頁。

\textbf{公式正答：} 3（single）。

\paragraph{問題の要旨} 風化・熱水変質・スレーキングが岩石強度へ及ぼす作用を判定する。

\Needspace{4\baselineskip}\paragraph{確認した条件・記述（独自要約）}\mbox{}\par

\noindent\textbf{方向：} 最も不適切なもの。\par

\noindent\textbf{対象：} 一般的な岩石分類・形成過程。\par

\paragraph{問われる核心} 風化・変質と間隙内圧による骨格破壊

\paragraph{必要知識} 風化・変質と間隙内圧による骨格破壊。母材・運搬堆積・有機物起源、風化、弱面、水の作用、岩盤と土の物性。

\paragraph{考え方・正答を選ぶ理由（査読待ち初稿）} スレーキングは水への浸漬などで岩塊が崩れる現象である。封入空気モデルでは、水の侵入で圧縮された空気の内圧が骨格を引っ張るので、岩塊の圧縮力とした3が不適切である。空気の状態変化と固体の応力を別々に記述する。膨潤・吸着・化学劣化なども機構となり得る。

\Needspace{4\baselineskip}\paragraph{選択肢の検討（査読待ち初稿）}\mbox{}\par

\noindent\textbf{1：適切（一般的分類）} 有色鉱物に富む黒い火成岩ほど一般に風化に不安定である。。高温形成のFe・Mgに富む鉱物は地表環境で変化しやすい。色そのものを原因とする法則ではなく、割れ目や組織も耐風化性を左右する。\par

\noindent\textbf{2：適切（一般的分類）} 磁鉄鉱から赤鉄鉱・褐鉄鉱への酸化が風化層の色に現れる。。Feの酸化と水和等により赤褐色系の鉱物を生じる。褐鉄鉱は単一の理想化学式を持つ一鉱物としてではなく含水酸化鉄の集合の名称として扱う。\par

\noindent\textbf{3：不適切} 浸水で間隙空気が圧縮され、岩塊内に圧縮力を生じてスレーキングする。。空気は圧縮されるが、内圧が孔壁・骨格を押し広げて引張破壊へ働く。空気の圧縮と岩塊への圧縮を混同している。封入空気だけが全試料の機構とも限らない。\par

\noindent\textbf{4：適切（一般的分類）} グリーンタフ地域の新第三紀層は熱水変質で軟弱化していることが多い。。火山性の岩石・堆積物が熱水反応で粘土化等を受け、低強度となる場合がある。全てのグリーンタフが同じ弱さという意味ではない。\par

\noindent\textbf{5：適切（一般的分類）} 一般に風化・変質は弱化させるが、珪化は強度を増すこともある。。シリカの充填・交代は緻密化や結合強化をもたらす場合があり、変質=常に強度低下とはしない。\par

\paragraph{再利用する考え方} 風化・変質と間隙内圧による骨格破壊を成因と観察される性質に結び付け、例示と全称命題、主要鉱物と副成分、公式解答と未解決の解釈を区別する。

\paragraph{共通理論} \hyperref[theory:engineering:weathering]{第\ref*{theory:engineering:weathering}節}, \hyperref[theory:engineering:slaking]{第\ref*{theory:engineering:slaking}節}

\par 参考：\cite{slaking1979,slaking2006,jseg_engineering_faq}。

\par 難易度 3（校準初稿）、優先度未評価。

\section{特殊土の母材・成因}

\subsection{平成25年度 問19}
\label{q:h25-reg-applied-19}

\noindent\textbf{原典：} sources/H25/17\_applied\_science.pdf、PDF物理8頁。

\textbf{公式正答：} 5（single）。

\paragraph{問題の要旨} まさ土・しらす・黄土・火山灰質粘性土・泥炭の成因について適切な記述を選ぶ。

\Needspace{4\baselineskip}\paragraph{確認した条件・記述（独自要約）}\mbox{}\par

\noindent\textbf{方向：} 最も適切なもの。\par

\noindent\textbf{対象：} 土の一般的な名称と成因。\par

\paragraph{問われる核心} 特殊土の母材・成因・残存植物組織

\paragraph{必要知識} 特殊土の母材・成因・残存植物組織。母材・運搬堆積・有機物起源、風化、弱面、水の作用、岩盤と土の物性。

\paragraph{考え方・正答を選ぶ理由（査読待ち初稿）} 色や粒状の見た目だけでなく母材と形成過程を対応させる。まさ土は花崗岩風化、しらすは火山噴出物、黄土は風成シルト、火山灰質粘性土は火山灰の風化等、泥炭は植物遺体の不完全分解である。適切なのは泥炭の5である。

\Needspace{4\baselineskip}\paragraph{選択肢の検討（査読待ち初稿）}\mbox{}\par

\noindent\textbf{1：不適切} まさ土の母材は流紋岩である。。まさ土は主に花崗岩類の風化により砂状となった土である。残留と崩積はその後の移動の区別であり、流紋岩を母材とする記述が誤り。\par

\noindent\textbf{2：不適切} しらすは花崗岩の風化物とその二次堆積物である。。しらすは主に珪長質火砕流の未溶結部などに由来する火山性堆積物で、二次堆積物もある。花崗岩の風化土とは起源が違う。\par

\noindent\textbf{3：不適切} 黄土は海洋堆積物である。。黄土（レス）は風で運ばれたシルトを主体とする風成堆積物で、海成という説明が誤り。\par

\noindent\textbf{4：不適切} 火山灰質粘性土をレスと呼ぶ。。火山灰由来の粘性土と風成シルトのレスを同一にしない。分布地域が一部重なるだけでは成因を一致させられない。\par

\noindent\textbf{5：適切} 泥炭はピートとも呼ばれ、未分解の植物組織を肉眼で判定できる。。水分の多い環境などで植物遺体の分解が遅れ、繊維・組織を残す有機質土が泥炭である。分解程度には幅があり、あらゆる泥炭の全組織が必ず肉眼で見えるという一般化は避ける。\par

\paragraph{再利用する考え方} 特殊土の母材・成因・残存植物組織を成因と観察される性質に結び付け、例示と全称命題、主要鉱物と副成分、公式解答と未解決の解釈を区別する。

\paragraph{共通理論} \hyperref[theory:engineering:soil]{第\ref*{theory:engineering:soil}節}

\par 参考：\cite{mlit_masa,maff_shiras2017,mlit_peat,gsj_loess2002}。

\paragraph{関連問題} \hyperref[q:h27-reg-applied-29]{H27-REG-APPLIED-29}

\par 難易度 3（校準初稿）、優先度未評価。

\subsection{平成27年度 問29}
\label{q:h27-reg-applied-29}

\noindent\textbf{原典：} sources/H27/17\_applied\_science.pdf、PDF物理11頁。

\textbf{公式正答：} 5（single）。

\paragraph{問題の要旨} まさ土・しらす・黄土・火山灰質粘性土・泥炭の成因について適切な記述を選ぶ。

\Needspace{4\baselineskip}\paragraph{確認した条件・記述（独自要約）}\mbox{}\par

\noindent\textbf{方向：} 最も適切なもの。\par

\noindent\textbf{対象：} 土の一般的な名称と成因。\par

\paragraph{問われる核心} 特殊土の母材・成因・残存植物組織

\paragraph{必要知識} 特殊土の母材・成因・残存植物組織。母材・運搬堆積・有機物起源、風化、弱面、水の作用、岩盤と土の物性。

\paragraph{考え方・正答を選ぶ理由（査読待ち初稿）} 色や粒状の見た目だけでなく母材と形成過程を対応させる。まさ土は花崗岩風化、しらすは火山噴出物、黄土は風成シルト、火山灰質粘性土は火山灰の風化等、泥炭は植物遺体の不完全分解である。適切なのは泥炭の5である。

\Needspace{4\baselineskip}\paragraph{選択肢の検討（査読待ち初稿）}\mbox{}\par

\noindent\textbf{1：不適切} まさ土の母材は流紋岩である。。まさ土は主に花崗岩類の風化により砂状となった土である。残留と崩積はその後の移動の区別であり、流紋岩を母材とする記述が誤り。\par

\noindent\textbf{2：不適切} しらすは花崗岩の風化物とその二次堆積物である。。しらすは主に珪長質火砕流の未溶結部などに由来する火山性堆積物で、二次堆積物もある。花崗岩の風化土とは起源が違う。\par

\noindent\textbf{3：不適切} 黄土は海洋堆積物である。。黄土（レス）は風で運ばれたシルトを主体とする風成堆積物で、海成という説明が誤り。\par

\noindent\textbf{4：不適切} 火山灰質粘性土をレスと呼ぶ。。火山灰由来の粘性土と風成シルトのレスを同一にしない。分布地域が一部重なるだけでは成因を一致させられない。\par

\noindent\textbf{5：適切} 泥炭はピートとも呼ばれ、未分解の植物組織を肉眼で判定できる。。水分の多い環境などで植物遺体の分解が遅れ、繊維・組織を残す有機質土が泥炭である。分解程度には幅があり、あらゆる泥炭の全組織が必ず肉眼で見えるという一般化は避ける。\par

\paragraph{再利用する考え方} 特殊土の母材・成因・残存植物組織を成因と観察される性質に結び付け、例示と全称命題、主要鉱物と副成分、公式解答と未解決の解釈を区別する。

\paragraph{共通理論} \hyperref[theory:engineering:soil]{第\ref*{theory:engineering:soil}節}

\par 参考：\cite{mlit_masa,maff_shiras2017,mlit_peat,gsj_loess2002}。

\paragraph{関連問題} \hyperref[q:h25-reg-applied-19]{H25-REG-APPLIED-19}, \hyperref[q:r01-ret-applied-28]{R01-RET-APPLIED-28}

\par 難易度 3（校準初稿）、優先度未評価。

\subsection{令和元年度再試験 問28}
\label{q:r01-ret-applied-28}

\noindent\textbf{原典：} sources/R01\_retest/17\_applied\_science.pdf、PDF物理9頁。

\textbf{公式正答：} 3（single）。

\paragraph{問題の要旨} まさ土・しらす・黄土・火山灰質粘性土・泥炭の成因について適切な記述を選ぶ。

\Needspace{4\baselineskip}\paragraph{確認した条件・記述（独自要約）}\mbox{}\par

\noindent\textbf{方向：} 最も適切なもの。\par

\noindent\textbf{対象：} 土の一般的な名称と成因。\par

\paragraph{問われる核心} 特殊土の母材・成因・残存植物組織

\paragraph{必要知識} 特殊土の母材・成因・残存植物組織。母材・運搬堆積・有機物起源、風化、弱面、水の作用、岩盤と土の物性。

\paragraph{考え方・正答を選ぶ理由（査読待ち初稿）} 色や粒状の見た目だけでなく母材と形成過程を対応させる。まさ土は花崗岩風化、しらすは火山噴出物、黄土は風成シルト、火山灰質粘性土は火山灰の風化等、泥炭は植物遺体の不完全分解である。適切なのは泥炭の3である。

\Needspace{4\baselineskip}\paragraph{選択肢の検討（査読待ち初稿）}\mbox{}\par

\noindent\textbf{1：不適切} 黄土は海洋堆積物である。。黄土（レス）は風で運ばれたシルトを主体とする風成堆積物で、海成という説明が誤り。\par

\noindent\textbf{2：不適切} しらすは花崗岩の風化物とその二次堆積物である。。しらすは主に珪長質火砕流の未溶結部などに由来する火山性堆積物で、二次堆積物もある。花崗岩の風化土とは起源が違う。\par

\noindent\textbf{3：適切} 泥炭はピートとも呼ばれ、未分解の植物組織を肉眼で判定できる。。水分の多い環境などで植物遺体の分解が遅れ、繊維・組織を残す有機質土が泥炭である。分解程度には幅があり、あらゆる泥炭の全組織が必ず肉眼で見えるという一般化は避ける。\par

\noindent\textbf{4：不適切} まさ土の母材は流紋岩である。。まさ土は主に花崗岩類の風化により砂状となった土である。残留と崩積はその後の移動の区別であり、流紋岩を母材とする記述が誤り。\par

\noindent\textbf{5：不適切} 火山灰質粘性土をレスと呼ぶ。。火山灰由来の粘性土と風成シルトのレスを同一にしない。分布地域が一部重なるだけでは成因を一致させられない。\par

\paragraph{再利用する考え方} 特殊土の母材・成因・残存植物組織を成因と観察される性質に結び付け、例示と全称命題、主要鉱物と副成分、公式解答と未解決の解釈を区別する。

\paragraph{共通理論} \hyperref[theory:engineering:soil]{第\ref*{theory:engineering:soil}節}

\par 参考：\cite{mlit_masa,maff_shiras2017,mlit_peat,gsj_loess2002}。

\paragraph{関連問題} \hyperref[q:h27-reg-applied-29]{H27-REG-APPLIED-29}

\par 難易度 3（校準初稿）、優先度未評価。

\section{未確定}

\subsection{平成28年度 問24}
\label{q:h28-reg-applied-24}

\noindent\textbf{原典：} sources/H28/17\_applied\_science.pdf、PDF物理10頁。

\textbf{公式正答：} 3（single）。

\par\textit{原典本文の校正と独自解説は未作成。分類のみの登録項目。}

\subsection{令和元年度 問29}
\label{q:r01-reg-applied-29}

\noindent\textbf{原典：} sources/R01/17\_applied\_science.pdf、PDF物理10頁。

\textbf{公式正答：} 5（single）。

\par\textit{原典本文の校正と独自解説は未作成。分類のみの登録項目。}

\subsection{令和6年度 問18}
\label{q:r06-reg-applied-18}

\noindent\textbf{原典：} sources/R06/17\_applied\_science.pdf、PDF物理6頁。

\textbf{公式正答：} 5（single）。

\par\textit{原典本文の校正と独自解説は未作成。分類のみの登録項目。}
